Compares memory with the expected (:term:base.terms.rn_register:) and conditionally stores the desired (:term:base.terms.rn_register:). After the operation, the expected register contains the observed old memory value. (:ref:base.registers.STATE.FLAGS.Z:)\texttt{=1} means the store happened; (:ref:base.registers.STATE.FLAGS.Z:)\texttt{=0} means memory was left unchanged.


The memory read, optional memory write, expected-register update, and flag update form one atomic read-modify-write operation.
The encoded memory-order selector governs the complete operation on both comparison outcomes. On failure, the
operation contributes the observed memory read as its ordering event: acquire orders later (:term:base.terms.memory_operation|plural:) after
that read, release orders earlier (:term:base.terms.memory_operation|plural:) before that read, and sequential consistency places that read in
the global SC order. On success, the memory write carries the same selected order to observers of that write. A
failed comparison contributes no memory write and creates no release sequence.

\begin{itemize}
\item If memory equals the original expected value, the desired value is stored and (:ref:base.registers.STATE.FLAGS.Z:) is set.
\item Otherwise memory is unchanged and (:ref:base.registers.STATE.FLAGS.Z:) is cleared.
\item In both cases the expected register receives the value originally read from memory, and (:ref:base.registers.STATE.FLAGS.N:), (:ref:base.registers.STATE.FLAGS.C:), and (:ref:base.registers.STATE.FLAGS.V:) are cleared.
\end{itemize}
A memory-access fault exposes none of these updates.
